






Bisher haben wir uns mit 








Ein Tangentialraum $T_x M$ mit Metrik kennt zunächst nur die quadratische Struktur
$$
v \otimes w \mapsto\langle v, w\rangle,
$$
also genau das, was in der Laplace-Geometrie auftaucht. Die Clifford-Algebra $\mathrm{Cl}\left(T_x M, g_x\right)$ ist die algebraische Hülle dieser Metrik, in der Vektoren nicht nur skalar multipliziert werden, sondern die Relation
$$
v \cdot w+w \cdot v=-2\langle v, w\rangle
$$
erfüllen. Ein Spinorraum $\Sigma_x$ ist dann ein irreduzibler Modul dieser Algebra. Der Dirac-Operator greift genau auf diese Struktur zu: Sein Prinzipalsymbol
$$
\sigma_1(D, \xi)=c(\xi)
$$
ist kein Skalar, sondern die Clifford-Multiplikation des Tangentialraums selbst.
Er "sieht" also nicht nur die Metrik, sondern die volle linearisierte Spin-Geometrie jedes Tangentialraums.